% !TeX program = xelatex
% =====================================================================
% ROZDZIAŁ FORMALNY — doktorat (język polski; publikacje po angielsku)
% „Warunki przyczynowości i wycieki danych w strumieniowej ewaluacji
%  detektorów na grafach komunikacji"
%
% Powiązanie z artykułami:
%   * artykuly_bogdan/A-Journal  (lateral phishing, detekcja kaskady)
%   * artykuly_bogdan/B-Conference (multipleks OSINT, wyciek rytmu)
% Każde twierdzenie mapuje się na eksperyment z paczki reprodukcyjnej
% (patrz tabela \ref{tab:mapping}).
% Kompilacja: tectonic formal_chapter.tex (XeTeX + polyglossia)
% =====================================================================
\documentclass[11pt]{article}
\usepackage{amsmath,amssymb,amsthm,amscd}
\usepackage{booktabs}
\usepackage{polyglossia}
\setmainlanguage{polish}
\usepackage{fontspec}
\usepackage[colorlinks]{hyperref}

\newtheorem{theorem}{Twierdzenie}
\newtheorem{proposition}[theorem]{Stwierdzenie}
\newtheorem{lemma}[theorem]{Lemat}
\newtheorem{corollary}[theorem]{Wniosek}
\newtheorem{definition}[theorem]{Definicja}
\newtheorem{remark}[theorem]{Uwaga}
\newtheorem{example}[theorem]{Przykład}

\newcommand{\E}{\mathbb{E}}
\newcommand{\Prob}{\mathbb{P}}
\newcommand{\R}{\mathbb{R}}
\newcommand{\event}{\mathbf{e}}
\newcommand{\strict}{\prec}

\title{Warunki przyczynowości i wycieki danych\\
w strumieniowej ewaluacji detektorów na grafach komunikacji}
\author{Rozdział formalny (robocza wersja rozprawy)}
\date{listopad 2026}

\begin{document}
\maketitle

\begin{abstract}
Rozdział aksjomatyzuje warunki ewaluacji, pod jakimi empiryczny AUC
detektora strumieniowego mówi cokolwiek o działaniu w warunkach
produkcyjnych, oraz stawia twierdzenia o tym, co się dzieje po ich
naruszeniu. Obiektem centralnym są: (i)~totalny porządek zdarzeń i klasa
\emph{zdarzeniowo-przyczynowych} cech i detektorów; (ii)~cztery
mechanizmy wycieku---look-ahead, cechy definujące etykietę, sygnatura
iniekcji i ślepota wewnątrzkubełkowa---każdy wyrażony jako twierdzenie
lub stwierdzenie z konstrukcyjnym kontrprzykładem;
(iii)~dekompozycja wyniku detektora na składową czysto czasową i
składową sprzęgającą czas z grafem, operacjonalizowana przez sondę
czysto czasową i kontrolę tasowania czasu; (iv)~wynik o
nierozróżnialności śladów przy realistycznej fabrykacji. Każdy obiekt
formalny jest instancjonowany zmierzonymi wielkościami rozdziałów
empirycznych.
\end{abstract}

\tableofcontents

% =====================================================================
\section{Ramka: strumienie zdarzeń i detektory}
\label{sec:setting}

\begin{definition}[Strumień zdarzeń]
\label{def:stream}
\emph{Strumień zdarzeń} to skończona sekwencja
$S=(\event_1,\dots,\event_N)$ zdarzeń
$\event=(s,v,b,\rho)$, gdzie $s$ to nadawca, $v$---odbiorca,
$b\in\{0,\dots,B-1\}$---\emph{kubełek czasowy}, a
$\rho\in\{1,\dots,N\}$---pozycja zdarzenia w ustalonym (zamrożonym)
porządku strumienia. Para
\[
t(\event) \;=\; (b,\rho)
\]
wyznacza \emph{porządek totalny} na zdarzeniach: $t(\event)\strict
t(\event')$ wtedy i tylko wtedy, gdy $b<b'$ lub $b=b'$ oraz
$\rho<\rho'$. Każde zdarzenie niesie etykietę $y(\event)\in\{0,1\}$
(legalne / atak).
\end{definition}

Pozycja $\rho$ to \emph{reguła remisów}: wewnątrz kubełka zdarzenia są
uporządkowane przez zamrożony porządek strumienia. W realnym logu
pocztowym porządek wewnątrz kubełka to porządek doręczeń; w strumieniu
eksperymentalnym jest ustalany raz na przebieg. Zamrożenie $\rho$
\emph{przed} liczeniem cech sprawia, że jednoczesność jest dobrze
zdefiniowana (sekcja~\ref{sec:withinbucket}).

\begin{definition}[Historia]
Dla zdarzenia $\event$ jego \emph{historia właściwa} to
\[
H_{\strict \event} \;=\; \{\event'\in S:\; t(\event')\strict t(\event)\}.
\]
\end{definition}

\begin{definition}[Cecha zdarzeniowo-przyczynowa]
\label{def:causal-feature}
Cecha $\varphi\colon S\times \{\event\}\to\R$ jest
\emph{zdarzeniowo-przyczynowa}, jeśli $\varphi(S,\event)$ jest mierzalną
funkcją $H_{\strict \event}$ i samego $\event$ oraz jest
\emph{niezmiennicza wobec usunięcia przyszłości}: dla każdego strumienia
$S'$ zgadzającego się z $S$ na $H_{\strict\event}\cup\{\event\}$ zachodzi
\[
\varphi(S,\event) \;=\; \varphi(S',\event).
\]
\end{definition}

Niezmienniczość wobec usunięcia przyszłości to operacyjny test z paczki
reprodukcyjnej (\texttt{tests/test\_causality.py}): cecha jest
przyczynowa wtedy i tylko wtedy, gdy usunięcie wszystkich zdarzeń o
$t\ge t(\event)$ nie zmienia jej wartości na $\event$.

\begin{definition}[Detektor przyczynowy]
\label{def:causal-detector}
\emph{Detektor} to funkcja $D\colon S\to\R^N$ przypisująca każdemu
zdarzeniu $\event_i$ wynik $D_i$. Detektor jest \emph{zdarzeniowo
przyczynowy}, jeśli jest mierzalną funkcją cech
zdarzeniowo-przyczynowych, a zbiór treningowy nie zawiera zdarzeń
późniejszych niż zdarzenia ewaluacyjne, na które może wpływać
(podział po ofiarach; sekcja~\ref{sec:split}).
\end{definition}

Empiryczną wielkością interesującą nas jest AUC detektora $D$ na
rozłącznym po ofiarach podzbiorze ewaluacyjnym $E\subseteq S$:
\[
\mathrm{AUC}(D,E) \;=\;
\Prob\big(D(\event) > D(\event') \mid \event\in E_1,\ \event'\in E_0\big)
\;+\;\tfrac12\,\Prob\big(D(\event)=D(\event')\big),
\]
gdzie $E_1,E_0$ to zdarzenia ataku i legalne z $E$.

% =====================================================================
\section{Co estymuje empiryczny AUC}
\label{sec:what-estimates}

\begin{definition}[Rozkład deplomentowy]
Rozkład generujący strumień $\mathcal{D}$ jest \emph{przyczynowy}, jeśli
samo generowanie danych respektuje porządek $t$: żadne zdarzenie nie
powstaje z użyciem informacji o zdarzeniach o większym $t$.
\end{definition}

\begin{theorem}[Obciążenie look-ahead]
\label{thm:lookahead}
Niech $\mathcal{D}$ będzie dowolnym przyczynowym rozkładem strumienia,
a $\varphi$ cechą \emph{niespełniającą} definicji
\ref{def:causal-feature}, tj.\ istnieje para strumieni $S,S'$
zgadzających się na $H_{\strict\event}\cup\{\event\}$, dla której
$\varphi(S,\event)\ne\varphi(S',\event)$. Wtedy dla każdego
$\varepsilon>0$ istnieją przyczynowe $\mathcal{D}$ i nieprzyczynowa
cecha $\varphi$ takie, że:
\begin{enumerate}
\item[(a)] każdy detektor zdarzeniowo-przyczynowy ma
      $\mathrm{AUC}\le \tfrac12+\varepsilon$ na $\mathcal{D}$
      (brak sygnału nadającego się do deplomentu), podczas gdy
\item[(b)] sam próg na $\varphi$ osiąga $\mathrm{AUC}\ge 1-\varepsilon$.
\end{enumerate}
Przy tym luka jest \emph{nieograniczona}: nie istnieje funkcja
$f\colon[0,1]\to[0,1]$ ograniczająca AUC „przeciekającej" cechy dane AUC
przyczynowej.
\end{theorem}

\begin{proof}[Dowód (konstrukcyjny; „bit rytmu")]
Niech strumień składa się ze zdarzeń legalnych, których kubełki są
losowane z \emph{pełnohistorycznego} zbioru aktywności ich nadawcy,
oraz ze zdarzeń ataku, których kubełki pochodzą z jego dopełnienia; obie
klasy mają te same rozkłady nadawców, odbiorców i brzegowe kubełków.
Zdefiniujmy $\varphi(\event)=\mathbf{1}[b \in A_s]$, gdzie $A_s$ to
\emph{pełnohistoryczny} zbiór aktywności nadawcy $s$, liczony z
całego $S$ wraz z okresem ewaluacji.
(a) Każdy detektor zdarzeniowo-przyczynowy widzi z konstrukcji dwie
wymienne podklasy kubełków: wewnątrz $H_{\strict\event}$ procesy
generowania klasy legalnej i ataku zbiegają, więc jego rozkład jest
niezmienniczy wobec przestawień etykiet niewidzianych zdarzeń; jego AUC
zbiega w oczekiwaniu do $\tfrac12$ (moc $tfrac12+\varepsilon$ osiągamy
biorąc dostatecznie dużo ziaren).
(b) Bit $\varphi$ jest z konstrukcji równy etykiecie, więc próg na nim
rozdziela klasy idealnie z dokładnością do błędu próbkowania
kontrolowanego przez $\varepsilon$.
Nieograniczoność: przeskalujmy konstrukcję prewalencją ataku
$\pi\to0$; przyczynowy AUC pozostaje $\tfrac12$, a przeciekający $1$.
\end{proof}

\begin{remark}
Twierdzenie~\ref{thm:lookahead} jest obiektem formalnym stojącym za
\emph{pułapką wycieku rytmu} z drugiego artykułu: estymacja rytmu
nadawcy na okresie ewaluacji daje pozorną czułość $0{,}73$, która
spada do zera podziałem czasowym (zmierzone: Tab.~IV artykułu B;
czułość pełnego multipleksu $0{,}51\to0{,}14\to0{,}01$ w miarę
zacieśniania kontroli).
\end{remark}

% =====================================================================
\section{Cechy definujące etykietę i warunek podziału}
\label{sec:split}

\begin{proposition}[Podział po ofiarach]
\label{prop:split}
Niech $D$ będzie trenowany na podzbiorze $S^{\mathrm{tr}}$ i
ewaluowany na $S^{\mathrm{te}}$. Jeśli jakaś ofiara należy do obu, to
dowolna cecha agregująca po zdarzeniach z otoczenia tej ofiary może
wykorzystać zdarzenia ataku z okresu ewaluacji w treningu
(\emph{samo-wyciek}), i twierdzenie~\ref{thm:lookahead} stosuje się
ponownie. Zatem podział ewaluacyjny musi być rozłączny po ofiarach.
\end{proposition}

\begin{corollary}[Wariant indukcyjny]
Jeśli dodatkowo \emph{organizacja} (spójna składowa grafu kontaktów)
jest rozłączna między treningiem a ewaluacją, estymujemy \emph{ryzyko
transferu} detektora, a nie jego dokładność w rozkładzie. Zmierzona
różnica między tymi protokołami jest sama w sobie wynikiem (cechy
wolumetryczne transferują się najlepiej; sygnał odbiorczy degraduje).
\end{corollary}

% =====================================================================
\section{Ślepota wewnątrzkubełkowa}
\label{sec:withinbucket}

Przyczynowość kubełkowa (okna ograniczone do kubełków $\le b-1$) to
warunek \emph{słabszy} niż przyczynowość zdarzeniowa: ujmuje detektorowi
prefiks kubełka $b$.

\begin{proposition}[Widoczna część burstu]
\label{prop:withinbucket}
Niech burst ataku umieszcza $m\ge2$ zdarzeń ataku w tym samym kubełku
$b$, dotykających wspólnego otoczenia ofiary. Detektor zdarzeniowo
przyczynowy z oknem $W$ kubełków widzi, dla $j$-tego zdarzenia burstu
($j=2,\dots,m$), dokładnie $j-1$ wcześniejszych zdarzeń ataku tego
burstu; detektor kubełkowy nie widzi żadnego. W konsekwencji oczekiwana
widoczna przeszła część burstu dla detektora kubełkowego jest zerowa
wewnątrz kubełka, a cały sygnał ko-występowania jest dla niego
niewidoczny, podczas gdy liczniki wolumetryczne---przyczynowe w sensie
definicji~\ref{def:causal-feature}---widzą burst narastający.
\end{proposition}

\begin{remark}
Zmierzone w reżimach blast na realnych grafach: odsetek zdarzeń ataku,
których fala otoczenia przypada \emph{ściśle w przeszłych kubełkach},
wynosi $0{,}21$--$0{,}25$ (Enron; \texttt{exp\_blast\_inversion.csv}),
podczas gdy $0{,}58$--$0{,}64$ ofiar ataku jest trafionych wielokrotnie
\emph{w tym samym kubełku}. Przy przyczynowości zdarzeniowej ta
asymetria jest mechanizmem inwersji blast: detektory odbiorcze przegrywają
($\mathrm{AUC}\approx0{,}62$) dokładnie tam, gdzie wewnątrzkubełkowe
liczniki wolumetryczne wygrywają ($\mathrm{AUC}\approx0{,}83$). Żadne
odwołanie do szczegółu implementacyjne nie jest potrzebne, gdy reguła
remisów $\rho$ jest jawna.
\end{remark}

% =====================================================================
\section{Sygnatura iniekcji i sonda czasowa}
\label{sec:injection}

Ataki wstrzykiwane niosą nie tylko wzorzec grafowy, ale też
\emph{brzegowy rozkład czasu}: rozkład $Q$ odstępów między wysyłkami
atakującego.

\begin{lemma}[Wykrywalność brzegowa]
\label{lem:probe}
Niech odstępy legalne (per nadawca) będą i.i.d.\ z $P$, a odstępy ataku
z $Q$. Wtedy dwuprobowy test KS $D_{KS}(P,Q)$ ogranicza z góry AUC
optymalnego czysto czasowego detektora \emph{brzegowego}, a detektor ten
osiąga
\[
\mathrm{AUC}_{\mathrm{probe}} \;=\; \tfrac12 + \Theta\big(D_{KS}(P,Q)\big).
\]
\end{lemma}

\begin{proof}[Zarys dowodu]
Optymalny test na pojedynczy próg wartości odstępu ma przewagę
ograniczoną przez odległość total variation rozkładów odstępów,
która z góry jest ograniczona przez $2D_{KS}$ (standardowa relacja
dwuprobowo-ksowa); $D_{KS}$ ogranicza ją z dołu przez próg w punkcie KS.
\end{proof}

\begin{theorem}[Usunięcie sygnatury]
\label{thm:injection}
Zastąpienie rozkładu odstępów ataku $Q$ przez empiryczny rozkład tła
$\widehat P$ (odstępy próbkowane z legalnych nadawców tego samego
harnessu) sprowadza składową brzegową do podłogi próbkowania:
\[
\E\big[\mathrm{AUC}_{\mathrm{probe}}\big] \;\le\; \tfrac12 +
O\big(D_{KS}(\widehat P, P)\big),
\qquad D_{KS}(\widehat P,P)\to 0 \text{ gdy } n\to\infty.
\]
W konsekwencji każda marża detektora, która \emph{przeżywa} zastąpienie,
jest przypisywalna cechom sprzęgającym czas ze strukturą grafu
(członom interakcyjnym), a nie samemu brzegowi czasowemu.
\end{theorem}

\begin{remark}
Przed zastąpieniem: $D_{KS}=0{,}46$--$0{,}79$,
$\mathrm{AUC}_{\mathrm{probe}}=0{,}70$--$0{,}76$ na trzech realnych
grafach. Po zastąpieniu: $D_{KS}=0{,}04$--$0{,}23$, sonda spada do
$0{,}54$--$0{,}63$; temporalny GNN utrzymuje marżę $+0{,}10$ nad
baseline wolumetrycznym (Enron $0{,}77$ vs $0{,}67$; CollegeMsg
$0{,}69$ vs $0{,}59$), podczas gdy ręczne cechy odbiorcze spadają do
poziomu wolumetrycznego ($0{,}61$--$0{,}71$ wobec $0{,}57$--$0{,}67$).
Zgodnie z twierdzeniem~\ref{thm:injection} przeżywająca marża jest
sygnałem interakcyjnym czytanym przez model nauczony; marża ręcznych
cech była z pierwszego rzędu sygnaturą iniekcji.
% TODO(checkpoint): potwierdzić marżę Enron w CI; anomalii tasowania
% (0,80 vs 0,77) nie wyjaśnia lemat 5 — rozstrzygnięcie po boot_ci.
\end{remark}

\begin{remark}[Zastrzeżenie: przesunięcie gęstości]
Empiryczne zastąpienie odstępów zmienia \emph{gęstość} czasową ataku,
gdy rozkład tła jest gęstszy niż deterministyczny spread, który
zastępuje (zmierzone: AUC temporalnego GNN \emph{wzrósł} z $0{,}70$ do
$0{,}81$ w harnessu syntetycznym). Zastąpienie jest więc kontrolą na
\emph{sygnaturę}, a nie monotonicznym pokrętłem „skrytszego ataku";
główna teza harnessu musi być sformułowana jako niezmienniczość
kontrastu w obu reżimach czasowania, nie jako pogorszenie.
\end{remark}

% =====================================================================
\section{Zależność od porządku: kontrola tasowania}
\label{sec:shuffle}

\begin{proposition}[Kontrola tasowania ogranicza zależność od porządku]
\label{prop:shuffle}
Niech $\sigma$ będzie losowym przestawieniem przypisań kubełków, a
$D^\sigma = D\circ\sigma$ detektorem przetrenowanym na tasowanym
strumieniu. Wtedy
\[
\Delta_{\mathrm{porz}} \;=\; \mathrm{AUC}(D) - \mathrm{AUC}(D^\sigma)
\]
jest \emph{ograniczeniem z góry} wkładu samego porządku czasowego:
tasowanie niszczy też ko-występowanie wewnątrz kubełków, więc
$\Delta_{\mathrm{porz}}$ nie identyfikuje czystej składowej porządku.
Detektor z $\Delta_{\mathrm{porz}}\approx0$ może być nadal czasowy, jeśli
jego sygnał mieszka w ko-występowaniu; detektor z
$\Delta_{\mathrm{porz}}>0$ jest wrażliwy na porządek.
\end{proposition}

\begin{remark}
Zmierzone: $\Delta_{\mathrm{porz}}\approx0{,}10$--$0{,}13$ dla
temporalnego GNN w harnessach syntetycznych (wrażliwowość na porządek
zgodnie z tezą). Na Enronie w reżimie odstępów empirycznych wariant
tasowany osiąga wynik \emph{wyższy} niż uporządkowany ($0{,}80$ vs
$0{,}77$): ostrzeżenie, że część marży na realnym grafie to
ko-występowanie lub brzeg stopnia, a nie porządek---uczciwa teza brzmi
„przeżywa kontrolę sygnatury", nie „wymaga porządku".
% TODO(checkpoint): CI z per-seed dla $\Delta_{\mathrm{porz}}$ na Enron/emp.
\end{remark}

% =====================================================================
\section{Realistyczna fabrykacja śladów}
\label{sec:fabrication}

Atakujący z drugiego artykułu może sfabrykować ślad OSINT. Niech $F$
będzie rozkładem \emph{prawdziwych} śladów (zbiorów warstw) legalnych
nowych nadawców.

\begin{proposition}[Cechy istnienia pokonane przez realistyczną fabrykację]
\label{prop:fab}
Jeśli atakujący próbuje każdy sfabrykowany ślad i.i.d.\ z $F$ (ta sama
mieszanka warstw, ta sama liczba krawędzi co u śladów prawdziwych) i
wsadza go jako realną krawędź multipleksu, wtedy:
(i) dowolna cecha zależna tylko od \emph{istnienia} lub \emph{liczby}
śladów ma ten sam rozkład na zdarzeniach sfabrykowanych i prawdziwych
(nierozróżnialność cech istnienia), więc jej moc rozdzielająca na
sfałszowanych atakach zbiega do jej mocy na legalnych nowych
nadawcach---czyli do zera;
(ii) kontrola permutacji ofiar (shuffle) jest niezależna od tempa
fabrykacji, bo wsadzone krawędzie są związane z prawdziwą ofiarą i
niewidoczne po permutacji;
(iii) tylko cechy \emph{spójności}---sprzężenia między warstwami i z
prawdziwym pokryciem nadawcy---mogą zachować moc, a ich marża jest
ograniczona przez ułamek warstw ataku \emph{konfliktujących} z prawdziwym
pokryciem nadawcy.
\end{proposition}

\begin{proof}[Zarys dowodu]
(i) to równość rozkładów cechy z konstrukcji próbkowania.
(ii) zachodzi, bo permutacja usuwa wsadzoną krawędź z lookupu
(nadawca, permutowana ofiara). (iii) wynika z dekompozycji indykatora
warstwy na część prawdziwą i wsadzoną; wsadzona warstwa wnosi do
rozdzielania wyłącznie przez interakcję z indykatorami prawdziwego
pokrycia.
\end{proof}

\begin{remark}
Przy fabrykatorze legacy (dowolna warstwa, honorowana per zdarzenie)
kontrola tasowania rosła z tempem fabrykacji ($0{,}51\to0{,}86$ przy
$100\%$)---artefakt naruszenia (ii). Przy realistycznym fabrykatorze ze
stwierdzenia~\ref{prop:fab} kontrola tasowania jest płaska na
$0{,}50\pm0{,}01$ w całym zakresie $0$--$100\%$, a AUC multipleksu spada
do poziomu samych kontaktów przy $50\%$ ($0{,}70$) z częściowym
odbiciem przy $100\%$ ($0{,}81$), zgodnie z (i)+(iii): przy pełnej
fabrykacji każdy atak niesie ślad wyglądający na prawdziwy i tylko
rezydualne sprzężenia rozdzielają.
\end{remark}

% =====================================================================
\section{Baza i ograniczenie triażu}
\label{sec:baserate}

\begin{proposition}[Arytmetyka bazy]
Dla detektora o czułości $r$ i odsetku fałszywych alarmów $\alpha$ przy
prewalencji $\pi$ (udziale ataków w strumieniu) precyzja wynosi
\[
\mathrm{prec}(\pi) \;=\;
\frac{r\pi}{r\pi+\alpha(1-\pi)},
\qquad
\mathrm{FP\ na\ wykrycie} \;=\; \frac{\alpha(1-\pi)}{r\pi},
\]
obie monotoniczne w $\pi$. Przy realistycznej prewalencji
$\pi\sim10^{-4}$ i punkcie pracy $(r,\alpha)=(0{,}10;0{,}01)$ każde
wykrycie niesie $\sim10^3$ fałszywych alarmów; klasa detektorów jest
narzędziem \emph{rankingowym}.
\end{proposition}

Wniosek używany w całej rozprawie: twierdzenia formułujemy w kategoriach
AUC / czułości z jawnym $\alpha$, nigdy jako „wykrywanie" bez bazy.

% =====================================================================
\section{Podsumowanie obiektów formalnych i ich instancji}
\label{sec:mapping}

\begin{table}[h]\centering\small
\begin{tabular}{llll}
\toprule
Obiekt formalny & Mechanizm & Eksperyment & Zmierzony efekt \\
\midrule
Tw.~\ref{thm:lookahead} & look-ahead & podział rytmu (B) & $0{,}73\to0$ czułości \\
Stw.~\ref{prop:split} & samo-wyciek & podział po ofiarach & --- (protokół) \\
Stw.~\ref{prop:withinbucket} & ślepota kubełkowa & inwersja blast & $0{,}25$ trafień wcześniejszych \\
Lem.~\ref{lem:probe} & brzeg czasowy & sonda KS & $0{,}70$--$0{,}76$ \\
Tw.~\ref{thm:injection} & usunięcie sygnatury & odstępy empiryczne & sonda $\to0{,}54$--$0{,}63$ \\
Stw.~\ref{prop:shuffle} & ograniczenie porządku & tasowanie czasu & $\Delta=0{,}10$--$0{,}13$ \\
Stw.~\ref{prop:fab} & realistyczna fabrykacja & sweep fabrykacji & shuffle płaski $0{,}50$ \\
Stw. (baza) & prewalencja & tab. bazy & $\sim10^3$ FP/wykrycie \\
\bottomrule
\end{tabular}
\caption{Mapowanie obiektów formalnych na zmierzone wielkości
(\texttt{code/results/*.csv} paczki reprodukcyjnej).}
\label{tab:mapping}
\end{table}

\paragraph{Co rozdział wnosi do tezy rozprawy.}
Twierdzenia przekształcają empiryczną dyscyplinę rozprawy w warunki
sprawdzalne: ewaluacja detektorów strumieniowych jest \emph{poprawna}
tylko wtedy, gdy (a) cechy przechodzą test usuwania
(def.~\ref{def:causal-feature}), (b) podział jest rozłączny po ofiarach
(stw.~\ref{prop:split}), (c) brzeg czasowy jest zmierzony i usunięty
(tw.~\ref{thm:injection}), (d) fabrykacja jest losowana z rozkładu
prawdziwych śladów (stw.~\ref{prop:fab}). Każdy z tych warunków był---i w
literaturze, i we wcześniejszych wersjach naszych eksperymentów---domoślnie
naruszany; każde naruszenie zawyżało raportowane liczby. Uczciwe
twierdzenia detekcyjne rozprawy to dokładnie te, które zostają po
(a)--(d).

\end{document}
